% !TeX program = pdflatex
% =============================================================================
% Standortbestimmung (Multiple Choice, unbenotet): s-t-Diagramm
% Kantonsschule KFR -- Physik
%
% Diagnostisches Instrument: wird zweimal ausgefuellt -- einmal VOR und
% einmal NACH der Doppellektion zum s-t-Diagramm (siehe
% s-t-diagramm-lektionsplanung.tex). Deckt die Kernbegriffe des
% Arbeitsblatts ab und testet gezielt die in lernziele.md dokumentierten
% Lernschwierigkeiten als Distraktoren (nicht erfundene Strohmaenner).
%
% Loesungen ein-/ausblenden wie beim Arbeitsblatt: \printanswers (Lehrer-
% exemplar) bzw. \noprintanswers (Schuelerexemplar, Standard) weiter unten.
% =============================================================================
\documentclass[11pt,a4paper]{exam}

\usepackage[T1]{fontenc}
\usepackage[utf8]{inputenc}
\usepackage[ngerman]{babel}
\usepackage{lmodern}
\usepackage[a4paper,margin=2.2cm,headheight=14pt]{geometry}
\usepackage{amsmath}
\usepackage{amssymb}
\usepackage{siunitx}
\sisetup{per-mode=fraction,fraction-command=\frac}
\usepackage{array,booktabs,tabularx,multirow}
\usepackage{enumitem}
\usepackage{graphicx}
\usepackage{xcolor}
\usepackage{colortbl}
\usepackage{tikz}
\usepackage{physics}
\usepackage[most]{tcolorbox}
\usepackage{lastpage}
\usepackage{needspace}

\usetikzlibrary{arrows.meta,calc,angles,quotes}

\usepackage{setspace}
\setlength{\parindent}{0pt}
\setlength{\parskip}{0.6em}
\setstretch{1.08}
\setlist[itemize]{itemsep=0.4em}

% -----------------------------
% Farben (Corporate-Farbschema Physik KFR, wie im Arbeitsblatt)
% -----------------------------
\definecolor{titleblue}{HTML}{183A6B}
\definecolor{accentblue}{HTML}{2E6F9E}
\definecolor{lightblue}{HTML}{EAF4FB}
\definecolor{verylightblue}{HTML}{F7FBFE}
\definecolor{lightgray}{HTML}{F4F6F8}
\definecolor{darkgray}{HTML}{4A4A4A}
\definecolor{accentorange}{HTML}{D9720C}
\definecolor{accentpurple}{HTML}{7B2D8E}

% -----------------------------
% Boxen (reduziert auf das fuer die Standortbestimmung Noetige)
% -----------------------------
\tcbset{before skip=10pt, after skip=10pt}
\newtcolorbox{titlebox}{
  enhanced, colback=titleblue, colframe=titleblue, boxrule=0pt, arc=3mm,
  left=3mm,right=3mm,top=3mm,bottom=3mm
}
\newtcolorbox{infobox}{
  enhanced, breakable, colback=lightblue, colframe=accentblue, boxrule=0.7pt,
  arc=2mm, left=6mm,right=6mm,top=3.5mm,bottom=3.5mm
}

\newlength{\aufgabenindent}
\settowidth{\aufgabenindent}{10.\hskip\labelsep}
\newtcolorbox{taskbox}[1][]{
  enhanced, breakable, colback=verylightblue, colframe=accentblue,
  title={#1}, fonttitle=\bfseries, boxrule=0.7pt, arc=2mm,
  enlarge left by=-\aufgabenindent, width=\textwidth,
  left=6mm,right=6mm,top=3.5mm,bottom=3.5mm
}

% --- Loesungen ein-/ausblenden -----------------------------------------------
\noprintanswers
% \printanswers

% --- Unbenotet: keine Punkte drucken ----------------------------------------
\pointformat{}

% --- Kopf-/Fusszeile ---------------------------------------------------------
\pagestyle{headandfoot}
\cfoot{\textcolor{darkgray}{Seite \thepage\ von \pageref{LastPage}}}

\newcommand{\Kopfzeile}[4]{%
  \begin{titlebox}
    {\color{white}\sffamily\bfseries\Large #4}
  \end{titlebox}
  \noindent\textbf{Klasse:}~#1 \hfill \textbf{Datum:}~#2 \hfill \textbf{Thema:}~#3
    \hfill \textbf{Lehrperson:}~Loris Keller
  \lhead{\textcolor{darkgray}{#3}}
  \rhead{\textcolor{darkgray}{#4}}
}

\newcommand{\Klasse}{4a}
\newcommand{\Datum}{21.08.2026}
\newcommand{\Thema}{s-t-Diagramm}
\newcommand{\ArbeitsblattTitel}{Standortbestimmung: s-t-Diagramm}

\begin{document}

\Kopfzeile{\Klasse}{\Datum}{\Thema}{\ArbeitsblattTitel}

\begin{infobox}
\textbf{Zu diesem Test}\\
Dieser Test wird \textbf{zweimal} ausgefuellt: einmal \textbf{vor} und einmal
\textbf{nach} der Doppellektion zum s-t-Diagramm. Er wird \textbf{nicht
benotet}. Kreuze einfach an, was du im Moment fuer richtig haeltst, auch
wenn du unsicher bist oder raten musst. Pro Frage ist jeweils genau eine
Antwort richtig.

\medskip
\textbf{Name:}~\rule{5cm}{0.4pt} \hfill
$\square$~vor der Doppellektion \hfill $\square$~nach der Doppellektion
\end{infobox}

\begin{questions}

\question Das folgende $s$-$t$-Diagramm zeigt die Bewegung eines Zugs.
\begin{center}
\begin{tikzpicture}[scale=0.75]
  \draw[gray!20] (0,0) grid[xstep=1,ystep=1] (6,5);
  \draw[->] (-0.3,0) -- (6.7,0) node[right] {\scriptsize $t$ in min};
  \draw[->] (0,-0.3) -- (0,5.5) node[above] {\scriptsize $s$ in km};
  \foreach \x in {1,2,3,4,5,6} \draw (\x,0.06) -- (\x,-0.06) node[below] {\scriptsize \x};
  \foreach \y in {1,2,3,4,5} \draw (0.06,\y) -- (-0.06,\y) node[left] {\scriptsize \y};
  \draw[thick,titleblue] (0,1) -- (6,4);
  \draw[dashed,gray] (3,0) -- (3,2.5) -- (0,2.5);
  \filldraw[accentblue] (3,2.5) circle (1.8pt);
\end{tikzpicture}
\end{center}
Wo befindet sich der Zug zum Zeitpunkt $t=\SI{3}{\minute}$?
\begin{choices}
  \choice Bei $s=\SI{3}{\kilo\meter}$.
  \CorrectChoice Bei $s=\SI{2.5}{\kilo\meter}$.
  \choice Bei $s=\SI{1}{\kilo\meter}$.
  \choice Bei $s=\SI{4}{\kilo\meter}$.
\end{choices}

\question Eine Person geht auf einer geraden Strasse. Ihre Bewegung ergibt
das folgende $s$-$t$-Diagramm.
\begin{center}
\begin{tikzpicture}[scale=0.75]
  \draw[->] (-0.3,0) -- (6.7,0) node[right] {\scriptsize $t$};
  \draw[->] (0,-0.3) -- (0,4.5) node[above] {\scriptsize $s$};
  \draw[thick,titleblue] (0,0.5) -- (2.5,3.5) -- (4,3.5) -- (6,1);
\end{tikzpicture}
\end{center}
Welche Beschreibung passt zu diesem Diagramm?
\begin{choices}
  \choice Die Person geht einen Hügel hinauf, macht oben eine Pause und
  geht dann auf der anderen Seite wieder hinunter.
  \CorrectChoice Die Person entfernt sich zuerst von ihrem Startpunkt,
  bleibt dann kurz stehen und bewegt sich anschliessend wieder in Richtung
  Startpunkt zurück.
  \choice Die Person bewegt sich die ganze Zeit mit konstanter
  Geschwindigkeit vorwärts.
  \choice Die Person steht die ganze Zeit still.
\end{choices}

\question Das folgende Diagramm zeigt die Bewegung eines Aufzugs.
\begin{center}
\begin{tikzpicture}[scale=0.75]
  \draw[gray!20] (0,0) grid[xstep=1,ystep=1] (7,4);
  \draw[->] (-0.3,0) -- (7.7,0) node[right] {\scriptsize $t$ in s};
  \draw[->] (0,-0.3) -- (0,4.5) node[above] {\scriptsize $s$ in m};
  \draw[thick,titleblue] (0,0.5) -- (2,3) -- (5,3) -- (7,1.5);
  \node[titleblue,above] at (3.5,3.05) {\scriptsize waagrechter Abschnitt};
\end{tikzpicture}
\end{center}
Was macht der Aufzug während des \emph{waagrechten} Abschnitts im
Diagramm?
\begin{choices}
  \choice Er bewegt sich besonders schnell.
  \choice Er bewegt sich mit konstanter, aber kleiner Geschwindigkeit.
  \CorrectChoice Er steht still (Geschwindigkeit $=\SI{0}{\meter\per\second}$).
  \choice Er ändert gerade die Richtung.
\end{choices}

\question Vier Schülerinnen und Schüler zeichnen je ein $s$-$t$-Diagramm.
Welches davon kann die Bewegung eines realen Objekts \emph{nicht}
darstellen?
\begin{center}
\begin{tabular}{@{}cccc@{}}
\begin{tikzpicture}[scale=0.5]
  \draw[->] (-0.2,0) -- (3.2,0) node[right] {\scriptsize $t$};
  \draw[->] (0,-0.2) -- (0,3.2) node[above] {\scriptsize $s$};
  \draw[thick,titleblue] (0,0.3) -- (3,2.7);
  \node[below] at (1.5,-0.7) {\scriptsize A};
\end{tikzpicture}
&
\begin{tikzpicture}[scale=0.5]
  \draw[->] (-0.2,0) -- (3.2,0) node[right] {\scriptsize $t$};
  \draw[->] (0,-0.2) -- (0,3.2) node[above] {\scriptsize $s$};
  \draw[thick,titleblue] (0,0.5) -- (1.3,2) -- (1.3,0.7) -- (3,2.5);
  \node[below] at (1.5,-0.7) {\scriptsize B};
\end{tikzpicture}
&
\begin{tikzpicture}[scale=0.5]
  \draw[->] (-0.2,0) -- (3.2,0) node[right] {\scriptsize $t$};
  \draw[->] (0,-0.2) -- (0,3.2) node[above] {\scriptsize $s$};
  \draw[thick,titleblue] (0,2.5) -- (1.5,1.5) -- (3,2.5);
  \node[below] at (1.5,-0.7) {\scriptsize C};
\end{tikzpicture}
&
\begin{tikzpicture}[scale=0.5]
  \draw[->] (-0.2,0) -- (3.2,0) node[right] {\scriptsize $t$};
  \draw[->] (0,-0.2) -- (0,3.2) node[above] {\scriptsize $s$};
  \draw[thick,titleblue] (0,1.5) -- (3,1.5);
  \node[below] at (1.5,-0.7) {\scriptsize D};
\end{tikzpicture}
\end{tabular}
\end{center}
\begin{choices}
  \choice Diagramm A
  \CorrectChoice Diagramm B
  \choice Diagramm C
  \choice Diagramm D
\end{choices}

\question Zwei Fahrzeuge fahren auf derselben Strasse; das Diagramm zeigt
ihre Bewegung.
\begin{center}
\begin{tikzpicture}[scale=0.75]
  \draw[gray!20] (0,0) grid[xstep=1,ystep=1] (6,5);
  \draw[->] (-0.3,0) -- (6.7,0) node[right] {\scriptsize $t$ in s};
  \draw[->] (0,-0.3) -- (0,5.5) node[above] {\scriptsize $s$ in m};
  \draw[thick,titleblue] (0,3.5) -- (6,4.7) node[right] {\scriptsize A};
  \draw[thick,accentorange] (0,0.5) -- (6,4.2) node[right] {\scriptsize B};
  \draw[dashed,gray] (3,0) -- (3,5);
  \filldraw[titleblue] (3,4.1) circle (1.8pt);
  \filldraw[accentorange] (3,2.35) circle (1.8pt);
\end{tikzpicture}
\end{center}
Zum Zeitpunkt der gestrichelten Linie: Welches Fahrzeug ist in \emph{diesem
Moment} schneller unterwegs?
\begin{choices}
  \choice Fahrzeug A, weil es sich zu diesem Zeitpunkt weiter oben
  befindet (grösserer Ort $s$).
  \CorrectChoice Fahrzeug B, weil seine Linie an dieser Stelle steiler
  verläuft.
  \choice Beide sind gleich schnell, weil beide Linien Geraden sind.
  \choice Das lässt sich aus dem Diagramm nicht bestimmen.
\end{choices}

\question Ein Ball befindet sich auf einer Rampe. Der Nullpunkt der
$s$-Achse liegt am unteren Ende der Rampe, positive Werte von $s$ liegen
rampaufwärts. Das folgende Diagramm zeigt, wie der Ball die Rampe
hinunterrollt.
\begin{center}
\begin{tikzpicture}[scale=0.75]
  \draw[gray!20] (0,0) grid[xstep=1,ystep=1] (6,4);
  \draw[->] (-0.3,0) -- (6.7,0) node[right] {\scriptsize $t$ in s};
  \draw[->] (0,-0.3) -- (0,4.5) node[above] {\scriptsize $s$ in m};
  \draw[thick,titleblue] (0,3.5) -- (6,0.5);
\end{tikzpicture}
\end{center}
Was bedeutet die \emph{negative} Steigung in diesem Diagramm?
\begin{choices}
  \choice Die Zeit läuft in diesem Abschnitt rückwärts.
  \CorrectChoice Der Ort $s$ des Balls nimmt mit der Zeit ab: Er bewegt
  sich in die negative Richtung.
  \choice Der Ball wird langsamer, bis er stehen bleibt.
  \choice Das Diagramm ist physikalisch nicht möglich.
\end{choices}

\question Zwei Läuferinnen bewegen sich auf einer geraden Bahn. Das
Diagramm zeigt ihre Bewegung.
\begin{center}
\begin{tikzpicture}[scale=0.75]
  \draw[gray!20] (0,0) grid[xstep=1,ystep=1] (6,4);
  \draw[->] (-0.3,0) -- (6.7,0) node[right] {\scriptsize $t$ in s};
  \draw[->] (0,-0.3) -- (0,4.5) node[above] {\scriptsize $s$ in m};
  \draw[thick,titleblue] (0,0.5) -- (6,3.5) node[right] {\scriptsize 1};
  \draw[thick,accentorange] (0,3.5) -- (6,0.5) node[right] {\scriptsize 2};
\end{tikzpicture}
\end{center}
Welche Aussage über die beiden Läuferinnen trifft zu?
\begin{choices}
  \choice Läuferin 1 ist schneller als Läuferin 2, weil ihre Linie nach
  oben zeigt.
  \CorrectChoice Beide sind gleich schnell unterwegs (gleicher Betrag der
  Steigung), bewegen sich aber in entgegengesetzte Richtungen.
  \choice Läuferin 2 bewegt sich gar nicht, weil ihre Linie fällt.
  \choice Es lässt sich nicht vergleichen, da die Linien unterschiedlich
  verlaufen.
\end{choices}

\question Das folgende $s$-$t$-Diagramm zeigt eine lineare Bewegung, die
nicht bei $s=0$ beginnt.
\begin{center}
\begin{tikzpicture}[scale=0.75]
  \draw[gray!20] (0,0) grid[xstep=1,ystep=1] (6,6);
  \draw[->] (-0.3,0) -- (6.7,0) node[right] {\scriptsize $t$ in s};
  \draw[->] (0,-0.3) -- (0,6.5) node[above] {\scriptsize $s$ in m};
  \foreach \x in {1,2,3,4,5,6} \draw (\x,0.06) -- (\x,-0.06) node[below] {\scriptsize \x};
  \foreach \y in {1,2,3,4,5,6} \draw (0.06,\y) -- (-0.06,\y) node[left] {\scriptsize \y};
  \draw[thick,titleblue] (1,2) -- (5,6);
  \draw[dashed,gray] (1,0) -- (1,2) -- (0,2);
  \draw[dashed,gray] (5,0) -- (5,6) -- (0,6);
  \filldraw[black] (1,2) circle (1.6pt) node[below right] {\scriptsize $(t_1,s_1)$};
  \filldraw[black] (5,6) circle (1.6pt) node[above left] {\scriptsize $(t_2,s_2)$};
\end{tikzpicture}
\end{center}
Wie berechnest du die Geschwindigkeit dieser Bewegung richtig?
\begin{choices}
  \choice $v = \dfrac{s_2}{t_2} = \dfrac{\SI{6}{\meter}}{\SI{5}{\second}} = \SI{1.2}{\meter\per\second}$
  \CorrectChoice $v = \dfrac{s_2-s_1}{t_2-t_1} = \dfrac{\SI{6}{\meter}-\SI{2}{\meter}}{\SI{5}{\second}-\SI{1}{\second}} = \SI{1}{\meter\per\second}$
  \choice $v = \dfrac{s_1}{t_1} = \dfrac{\SI{2}{\meter}}{\SI{1}{\second}} = \SI{2}{\meter\per\second}$
  \choice $v = s_2 - s_1 = \SI{4}{\meter}$
\end{choices}

\question Ein Velofahrer startet bei $s=\SI{0}{\meter}$, fährt
$\SI{80}{\meter}$ in positive Richtung, kehrt danach um und fährt
$\SI{30}{\meter}$ in negative Richtung zurück.
\begin{center}
\begin{tikzpicture}[scale=0.9]
  \draw[->] (-0.3,0) -- (6,0) node[right] {\scriptsize $t$};
  \draw[->] (0,-0.3) -- (0,4.7) node[above] {\scriptsize $s$ in m};
  \draw[dashed,gray] (0,4) -- (3,4);
  \draw[dashed,gray] (0,2.5) -- (5,2.5);
  \node[left] at (0,4) {\scriptsize 80};
  \node[left] at (0,2.5) {\scriptsize 50};
  \node[left] at (0,0) {\scriptsize 0};
  \draw[thick,titleblue] (0,0) -- (3,4) -- (5,2.5);
  \filldraw[black] (0,0) circle (1.6pt);
  \filldraw[black] (3,4) circle (1.6pt);
  \filldraw[black] (5,2.5) circle (1.6pt);
\end{tikzpicture}
\end{center}
\begin{choices}
  \choice Der zurückgelegte Weg beträgt $\SI{50}{\meter}$.
  \CorrectChoice Der zurückgelegte Weg beträgt $\SI{110}{\meter}$.
  \choice Der zurückgelegte Weg ist gleich der Ortsänderung, also
  $\SI{50}{\meter}$.
  \choice Der zurückgelegte Weg lässt sich ohne Diagramm nicht bestimmen.
\end{choices}

\question Kann man einem Objekt zu einem einzelnen Zeitpunkt (z.\,B.\
genau bei $t=\SI{3}{\second}$) überhaupt eine Geschwindigkeit zuordnen?
\begin{center}
\begin{tikzpicture}[scale=0.8]
  \draw[->] (-0.3,0) -- (5.5,0) node[right] {\scriptsize $t$};
  \draw[->] (0,-0.3) -- (0,4.5) node[above] {\scriptsize $s$};
  \draw[thick,titleblue,domain=0:5,samples=50,smooth] plot (\x,{0.16*\x*\x});
  \draw[thick,accentpurple,dashed] (1.5,-0.06) -- (4.8,3.24) node[right] {\scriptsize Tangente};
  \filldraw[black] (3,1.44) circle (1.8pt) node[below right] {\scriptsize $t=\SI{3}{\second}$};
\end{tikzpicture}
\end{center}
\begin{choices}
  \choice Nein, für eine Geschwindigkeit braucht man immer zwei
  verschiedene, klar auseinanderliegende Messzeitpunkte. Zu einem
  einzelnen Zeitpunkt ist das unmöglich.
  \CorrectChoice Ja, das ist die Momentangeschwindigkeit. Sie entspricht
  der Steigung der Tangente an das $s$-$t$-Diagramm an diesem Punkt.
  \choice Nein, das ergibt nur bei einer gleichförmigen Bewegung Sinn.
  \choice Ja, aber nur, wenn sich das Objekt in diesem Moment gerade nicht
  bewegt.
\end{choices}

\question Welches der folgenden vier Diagramme zeigt eine
\textbf{\emph{gleichförmige}} Bewegung, bei der sich das Objekt tatsächlich
bewegt (konstante Geschwindigkeit ungleich null)?
\begin{center}
\begin{tabular}{@{}cccc@{}}
\begin{tikzpicture}[scale=0.5]
  \draw[->] (-0.2,0) -- (3.2,0) node[right] {\scriptsize $t$};
  \draw[->] (0,-0.2) -- (0,3.2) node[above] {\scriptsize $s$};
  \draw[thick,titleblue] (0,0.3) -- (3,2.7);
  \node[below] at (1.5,-0.7) {\scriptsize A};
\end{tikzpicture}
&
\begin{tikzpicture}[scale=0.5]
  \draw[->] (-0.2,0) -- (3.2,0) node[right] {\scriptsize $t$};
  \draw[->] (0,-0.2) -- (0,3.2) node[above] {\scriptsize $s$};
  \draw[thick,titleblue,domain=0:3,samples=40,smooth] plot (\x,{0.3*\x*\x});
  \node[below] at (1.5,-0.7) {\scriptsize B};
\end{tikzpicture}
&
\begin{tikzpicture}[scale=0.5]
  \draw[->] (-0.2,0) -- (3.2,0) node[right] {\scriptsize $t$};
  \draw[->] (0,-0.2) -- (0,3.2) node[above] {\scriptsize $s$};
  \draw[thick,titleblue] (0,1.5) -- (3,1.5);
  \node[below] at (1.5,-0.7) {\scriptsize C};
\end{tikzpicture}
&
\begin{tikzpicture}[scale=0.5]
  \draw[->] (-0.2,0) -- (3.2,0) node[right] {\scriptsize $t$};
  \draw[->] (0,-0.2) -- (0,3.2) node[above] {\scriptsize $s$};
  \draw[thick,titleblue] (0,0.3) -- (1.5,2.5) -- (3,1);
  \node[below] at (1.5,-0.7) {\scriptsize D};
\end{tikzpicture}
\end{tabular}
\end{center}
\begin{choices}
  \CorrectChoice Diagramm A
  \choice Diagramm B
  \choice Diagramm C
  \choice Diagramm D
\end{choices}

\question Eine Katze sitzt zunächst still, rennt dann für kurze Zeit sehr
schnell nach vorne und bleibt anschliessend wieder abrupt stehen. Welches
Diagramm passt zu dieser Beschreibung?
\begin{center}
\begin{tabular}{@{}ccc@{}}
\begin{tikzpicture}[scale=0.6]
  \draw[->] (-0.2,0) -- (4.2,0) node[right] {\scriptsize $t$};
  \draw[->] (0,-0.2) -- (0,3.2) node[above] {\scriptsize $s$};
  \draw[thick,titleblue] (0,0.4) -- (1,0.4) -- (2.5,2.8) -- (4,2.8);
  \node[below] at (2,-0.7) {\scriptsize A};
\end{tikzpicture}
&
\begin{tikzpicture}[scale=0.6]
  \draw[->] (-0.2,0) -- (4.2,0) node[right] {\scriptsize $t$};
  \draw[->] (0,-0.2) -- (0,3.2) node[above] {\scriptsize $s$};
  \draw[thick,titleblue] (0,0.4) -- (4,2.8);
  \node[below] at (2,-0.7) {\scriptsize B};
\end{tikzpicture}
&
\begin{tikzpicture}[scale=0.6]
  \draw[->] (-0.2,0) -- (4.2,0) node[right] {\scriptsize $t$};
  \draw[->] (0,-0.2) -- (0,3.2) node[above] {\scriptsize $s$};
  \draw[thick,titleblue] (0,2.8) -- (1,2.8) -- (2.5,0.4) -- (4,0.4);
  \node[below] at (2,-0.7) {\scriptsize C};
\end{tikzpicture}
\end{tabular}
\end{center}
\begin{choices}
  \CorrectChoice Diagramm A
  \choice Diagramm B
  \choice Diagramm C
\end{choices}

\end{questions}

\end{document}
